\documentclass{article} % For LaTeX2e
\usepackage{icomp2026_conference,times}

% Optional math commands from https://github.com/goodfeli/dlbook_notation.
\input{math_commands.tex}

\usepackage{hyperref}
\usepackage{url}
\usepackage{amsmath,amssymb,amsthm}
\usepackage{graphicx}
\usepackage{booktabs}

% ---------------------------------------------------------------------------
% TITLE  (working — refine before submission)
% ---------------------------------------------------------------------------
\title{Two Channels of Degradation in Recommender Feedback Loops:\\
Universal Taste Contraction and the Covariance Signature of Retraining}

% Authors must NOT appear in the submitted version. Keep \icompfinalcopy
% commented out until camera-ready. Non-anonymous submissions are rejected.
\author{Anonymous Authors}

%\icompfinalcopy % Uncomment ONLY for camera-ready, NOT for submission.

\begin{document}

\maketitle

\begin{abstract}
Industrial recommender systems are periodically retrained on the interaction
logs they themselves generate, closing a hidden feedback loop in which clicks
cease to be a neutral measurement of preference and observed quality metrics can
rise even as true utility falls. We study the resulting degradation in an
agent-based simulation and on real data, and show that it decomposes into two
distinct channels. The first is a universal $\beta$-channel: every click drags a
user's taste toward the consumed item, geometrically contracting diversity
regardless of whether the model is retrained. We prove this contraction
(Theorem~1) and confirm it across serving regimes and on MovieLens-20M, where it
accounts for roughly a $97\%$ loss of taste variance. The second is a
retraining-specific loop channel, and its signature is subtle: we find that the
commonly tracked collapse of total taste variance $\operatorname{tr}\hat\Sigma$
is driven by the $\beta$-channel, \emph{not} by retraining, whereas the loop's
own fingerprint lives in the \emph{shape} of the taste covariance and is captured
by the dispersion term of $\mathrm{KL}(P_T\Vert P_0)$ rather than by the trace.
This signature appears only when the model distinguishes users, the audience is
multi-modal, and the population is open; we further test the natural hypothesis
that retraining narrows the served distribution and refute it by direct
measurement, pointing instead to retargeting of drifted users. The results yield
a practical takeaway: the loop must be diagnosed through the observed--true
metric gap and the geometry of the taste covariance, not through headline
metrics, which improve precisely as the system degrades.
\end{abstract}

\section{Introduction}
% TODO

\section{Related work}
% TODO

\section{Problem setup}
% TODO

\section{Method / Theory}
% TODO

\section{Experiments}
% TODO

\section{Conclusion}
% TODO

\bibliography{references}
\bibliographystyle{icomp2026_conference}

\appendix
\section{Appendix}
% TODO

\end{document}
